\cotefile{157-2}
\scanpage{275}
\mspage{16}

3. $\mathbb{\Delta}_2^+,\mathbb{\Delta}_2^-$, correspondent aux diagrammes
\diagram{\begin{tikzcd}[ampersand replacement=\&] a_0 \arrow[r] \& a_1 \& \unsure{\text{---}|\longrightarrow} \& a_2 \end{tikzcd}}
resp.\quad \deleted{\unsure{$a$}}
\diagram{\begin{tikzcd}[ampersand replacement=\&] a_0 \& a_1 \arrow[l] \& a_2 \arrow[l] \end{tikzcd}},
on peut faire une remarque similaire que pour $\mathbb{\Delta}_1^+,\mathbb{\Delta}_1^-$

4\unsure{,|.} $\Phi_1,\Psi_1$, ce sont respectivement les intervalles
\diagram{\begin{tikzcd}[ampersand replacement=\&,row sep=2pt,column sep=small]
{} \& {} \& {} \& {} \& {} \& {} \& 0 \\
\Phi_1\colon \& 0 \& 1 \arrow[l,"{\underline{u}_0}"'] \arrow[r,"{\underline{u}_1}"] \& 2 \& \text{ou} \& 1 \arrow[ur] \arrow[dr] \& {} \\
{} \& {} \& {} \& {} \& {} \& {} \& 2
\end{tikzcd}}
\diagram{\begin{tikzcd}[ampersand replacement=\&,row sep=2pt,column sep=small]
{} \& {} \& {} \& {} \& {} \& {} \& 0 \arrow[dl] \\
\Psi_1\colon \& 0 \arrow[r,"{\underline{u}_0}"] \& 1 \& 2 \arrow[l,"{\underline{u}_1}"'] \& \unsure{\text{ou}} \& 1 \& {} \\
{} \& {} \& {} \& {} \& {} \& {} \& 2 \arrow[ul]
\end{tikzcd}}

\deleted{Cette fois $\Phi$} Ils correspondent respectivement aux diagrammes
\diagram{\begin{tikzcd}[ampersand replacement=\&] a_0 \& a_1 \arrow[l,"{u_0}"'] \arrow[r,"{u_1}"] \& a_2 \end{tikzcd}}
et
\diagram{\begin{tikzcd}[ampersand replacement=\&] a_0 \arrow[r] \& a_1 \& a_2 \arrow[l] \end{tikzcd}}

\unsure{Les|les} catégories $\Phi_1,\Psi_1$ sont \deleted{\unsure{évidemment}} opposées l'une de l'autre, et ne sont pas isomorphes. Aussi n'y a-t-il \inserted{en général} aucun iso. entre $\underline{\mathrm{Ch}}_{\underline{\Phi}_1}(X)$ et $\underline{\mathrm{Ch}}_{\underline{\Psi}_1}(X)$,

Soit $I$ un intervalle orienté, et considérons l'\uline{intervalle opposé} $I^{\mathrm{opp}}$, ayant $=$ catégorie sous-jacente mais l'orientation opposée. \deleted{Alors} Comme la définition de la \deleted{\unsure{de}} catégorie des chemins de type $I$ ne fait pas intervenir l'orientation de $I$, on a donc un iso. canonique
\[ \underline{\mathrm{Ch}}_I(X) \unsure{\xrightarrow{\sim}|\simeq} \underline{\mathrm{Ch}}_{I^{\mathrm{opp}}}(X), \]
% archive sheet number: 274
\scanpage{276}
\mspage{17}

Néanmoins on se gardera de \deleted{\illegible{1}} confondre un chemin \inserted{de type $I$} avec le chemin opposé de type $I^{\mathrm{opp}}$, la source de l'un \unsure{étant} le but de l'autre et inversement. C'est dire qu'un chemin doit être vu comme \uline{plus} qu'un \inserted{simple} foncteur
\diagram{\begin{tikzcd}[ampersand replacement=\&] I \arrow[r] \& X \end{tikzcd}},
il faut de plus se donner l'\uline{orientation de} $I$ comme faisant partie de la structure de chemin.

Donc la définition formelle raisonnable, qui colle à l'intuition \unsure{du chemin|des chemins} comme allant d'une \uline{source} bien définie vers un \uline{but} également bien défini, c'est qu'un chemin (défini en termes d'une cat. d'intervalle $I$) est la donnée d'un \uline{couple} $(\omega,c)$, où $\omega$ est une orientation de $I$, et \deleted{\unsure{$c$|,}}$c$ le diagramme $c\colon I\mathrel{\unsure{\text{---}|\longrightarrow}}X$.

Quand on regarde la catégorie des chemins d'un ``type'' déterminé, ce \unsure{type} est $(I,\omega)$, et non $I$. C'est quand $\omega$ est fixé \inserted{une bonne fois} \unsure{que} \deleted{la \unsure{ch}} cette catégorie s'identifie \unsure{simplement|impunément} à \deleted{\unsure{$\mathrm{Ch}$}} $\underline{\mathrm{Hom}}(I,X)$ sans plus.

Il faut voir cette catégorie de chemins comme munie de la structure essentielle
\diagram{\begin{tikzcd}[ampersand replacement=\&]
\underline{\mathrm{Ch}}_{(I,\omega)}\deleted{\unsure{(}}(X) \arrow[d,"{(\unsure{\sigma|\alpha},\beta)}"]\\
X\times X
\end{tikzcd}}
\[ \begin{aligned} c&\longmapsto(\sigma(c),\beta(c))\\[-2pt] &\text{où }\sigma(c)\text{ est la \uline{source},}\\ &\phantom{\text{où }}\beta(c)\text{ le but de }c. \end{aligned} \]
% archive sheet number: 275
\scanpage{277}
\mspage{18}

Ce foncteur s'identifie au foncteur de composition
\[ c\longmapsto c\circ\alpha \]
où $\alpha$ est le foncteur d'inclusion
\[ \left\{\begin{aligned}
\varepsilon&\xrightarrow{\alpha} I\\
\varepsilon&=\{0,1\},\quad\text{catégorie discrète à deux objets}\\
&\alpha\ \text{donné par}\quad \alpha(0)=\text{source de }I\\
&\phantom{\alpha\ \text{donné par}\quad}\alpha(1)=\text{but de }I
\end{aligned}\right. \]

\unsure{$\rightarrow$|,} Si $\omega'$ est l'\deleted{\unsure{caté}}orientation opposée de $I$, on a un diagramme commutatif
\diagram{\begin{tikzcd}[ampersand replacement=\&,row sep=large,column sep=large]
\underline{\mathrm{Ch}}_{(I,\omega)}(X) \arrow[r,dash,"{\varphi}"{yshift=4pt},"\sim"] \arrow[d,"{(\unsure{\sigma|\alpha},\beta)}"] \& \underline{\mathrm{Ch}}_{(I,\omega')}(X) \arrow[d,dash,"{(\sigma',\beta')}"]\\
X\times X \arrow[r,"s"] \& X\times X
\end{tikzcd}}
% slanted
\margin{Le foncteur $\varphi$ est celui qui associe à tout chemin $c$, de type $\tau$, le \uline{chemin opposé} $c^{\mathrm{opp}}$ ou $c^{-1}$, de type $\tau^{\mathrm{opp}}$. Si $\tau=\tau^{\mathrm{opp}}$, c'est donc un chemin \uline{\unsure{du|de} même type} $\tau$, \deleted{\unsure{l'aut}} $\varphi$ est un \uline{\unsure{autom}} de $\underline{\mathrm{Ch}}_{\tau}(X)_{\unsure{s}}$. Cet \unsure{autom} \inserted{(\unsure{naturel en|variant en} $X$)} est déduit d'un \unsure{autom.} $\theta$ de $I$, cf ci-dessous.}
$s$ est le foncteur de symétrie
\[ s(x,y)=(y,x). \]

Donc l'isomorphisme $\varphi$ ne respecte pas les flèches structurales ``source-but'' $(\sigma,\beta)$.

Il y a cependant un \unsure{cas heureux|mieux} quand il existe un iso. d'intervalles \uline{orientés}
\[ \theta\colon(I,\omega)\xrightarrow{\sim}(I,\omega'), \]
i.e. quand $I$ est \uline{symétrique}. Ce n'est pas le cas pour les $\mathbb{\Delta}_i$ $(i\geq 1)$ (ces catégories n'ont pas d'autre \unsure{autom.} que l'identité), mais c'est le cas pour les catégories $\Phi_1,\Psi_1$ de l'exemple $4^{\mathrm{o}}$, plus généralement \deleted{pour les}
% archive sheet number: 276
\scanpage{278}
\mspage{19}

pour les catégories réduites de longueur \uline{\unsure{pair|paire}}.

Comme chaque fois qu'on a un iso. $\theta$ d'int. \uline{orientés}, on trouve \deleted{cette fois} un iso. entre les catégories de chemins de type \deleted{$I$ et $I^{\mathrm{opp}}$} $(I,\omega)$ et $(I,\omega)^{\mathrm{opp}}=(I,\omega')$, \uline{compatible} cette fois avec les foncteurs $(\unsure{\sigma|\alpha},\beta)$ :
\diagram{\begin{tikzcd}[ampersand replacement=\&,column sep=small]
\underline{\mathrm{Ch}}_{(I,\omega)}(X) \arrow[rr,"{\substack{\unsure{\psi|\Psi}_{\theta}\\\sim}}"] \arrow[dr,"{(\sigma,\beta)}"] \&{} \& \underline{\mathrm{Ch}}_{(I,\omega')}(X) \arrow[dl,dash,"{(\sigma',\beta')}"']\\
{} \& X\times X \& {}
\end{tikzcd}}

\deleted{On peut regarder $\theta$ comme un \unsure{autom|autom.} non trivial de $I$ (échangeant \unsure{ses|les} deux extrémités) et qui opère \inserted{donc aussi} sur \deleted{\unsure{$\mathrm{Ch}$}} $\underline{\mathrm{Hom}}(I,X)$, et c'est ça \deleted{\unsure{la flèche}} l'iso. $\psi_\theta$ du \deleted{\unsure{le f}} diagramme précédent. Il transforme donc le \deleted{diagramme}\inserted{chemin de} type $\Phi_1$
\diagram{\begin{tikzcd}[ampersand replacement=\&,column sep=small,row sep=2pt]
{} \&{} \&{} \&{} \&{} \& a_0\\
a_0 \& a_1 \arrow[l] \arrow[r] \& a_2 \& \text{ou} \& a_1 \arrow[ur] \arrow[dr] \&{}\\
{} \&{} \&{} \&{} \&{} \&a_2
\end{tikzcd}}
\inserted{non pas en} en le chemin
\diagram{\begin{tikzcd}[ampersand replacement=\&,column sep=small,row sep=2pt]
{} \&{} \&{} \&{} \&{} \&a_2\\
a_2 \&a_1 \arrow[l] \arrow[r] \&a_0 \& \text{ou} \& a_1 \arrow[ur] \arrow[dr] \&{}\\
{} \&{} \&{} \&{} \&{} \&a_0
\end{tikzcd}}
\deleted{\unsure{Ils sont}} mais en lui-même. En tant que foncteur \deleted{\uline{Ces deux chemins} sont de même type $\Phi_1$, ($(I,\omega)\simeq(I,\omega')$) j'ai prétendu qu'il revient} \deleted{\unsure{à trouver|à trouve}} \inserted{sur $I$,} le chemin opposé à}
% archive sheet number: 277
\scanpage{279}
\mspage{20}

Il est bon de regarder \deleted{\unsure{les}} le \uline{type} de chemins comme une donnée \deleted{\unsure{d'un}} non catégorique, mais combinatoire, une simple suite $\tau$ de $n$ signes \inserted{$\tau_i$, $(1\leq i\leq n)$} $+$ ou $-$ ($n={}$ longueur des chemins envisagés). \deleted{\unsure{Et}} Cette donnée combinatoire définit un intervalle orienté \unsure{standard|standart}, \deleted{qui \unsure{pa}} $I_\tau$.

Quand on a un intervalle orienté quelconque $(I,\omega)$, il définit un type $\tau$, et on l'identifie canoniquement au type $I_\tau$. Il \deleted{\unsure{vaut}} vaut mieux écrire $\underline{\mathrm{Ch}}_\tau(X)$, \unsure{plus tôt|plutôt} que $\underline{\mathrm{Ch}}_{\unsure{(}I,\omega)}(X)$, pour bien indiquer que pour $X$ donné, c'est là une catégorie totalement ``épinglée'', \unsure{sans|des} automorphismes liés à la nature particulière de $\tau$ (\unsure{ô|à} les \unsure{autom.} de $\tau$ !).

On peut préciser ce point ainsi. Soit $\tau'$ un autre type d'intervalles, et proposons-nous de déterminer des \deleted{\unsure{fonct}} ``lois'' qui permettent, à tout \deleted{\unsure{che}} chemin de type $\tau$ dans une cat. \inserted{arbitraire} $X$, d'associer un chemin de type $\tau'$ dans cette \unsure{$\overline{\mathrm m}$|$=$} catégorie, et ceci de façon \inserted{``fonctorielle'', i.e.} compatible avec les \deleted{foncteurs} applications horizontales
\diagram{\begin{tikzcd}[ampersand replacement=\&,row sep=small]\mathrm{Ch}_\tau(X) \arrow[r,"{\mathrm{Ch}_\tau(f)}"] \arrow[d,"{L_X}"'] \& \mathrm{Ch}_\tau(X') \arrow[d,"{L_{X'}}"]\\
\mathrm{Ch}_{\tau'}(X) \arrow[r,"{\mathrm{Ch}_{\tau'}(f)}"] \&\mathrm{Ch}_{\tau'}(X')
\end{tikzcd}}
associées \deleted{aux foncteurs \unsure{q|p}} \inserted{à un} foncteur $f\colon X\mathrel{\unsure{\text{---}|\longrightarrow}}X'$.
% archive sheet number: 278
\scanpage{280}
\mspage{21}

Par \unsure{exemple|exple}, \unsure{prenant|prenons} le \deleted{chemin} type $\mathbb{\Delta}_2^+$ (i.e. $+1,+1$)
\diagram{\begin{tikzcd}[ampersand replacement=\&] 0 \arrow[r] \& 1 \arrow[r] \& 2 \end{tikzcd}},
correspondant aux chemins
\diagram{\begin{tikzcd}[ampersand replacement=\&] a_0 \arrow[r,dash,"{u_1}"] \& a_1 \arrow[r,"{u_2}"] \& a_2 \end{tikzcd}}
on peut associer à un tel chemin le chemin de type
\diagram{\begin{tikzcd}[ampersand replacement=\&] a_0 \arrow[r,"{u_2u_1}"] \& a_2 \end{tikzcd}},
ou aussi le chemin de type $\mathbb{\Delta}_2^+$
\diagram{\begin{tikzcd}[ampersand replacement=\&] a_0 \arrow[r,dash,"{u_2u_1}"] \& a_2 \arrow[r,"{\mathrm{id}}"] \& a_2 \end{tikzcd}},
\deleted{ou aussi} \deleted{\diagram{\begin{tikzcd}[ampersand replacement=\&] \unsure{a_0} \arrow[r,"{\mathrm{id}}"] \& a_0 \arrow[r,dash] \& \unsure{a_1} \arrow[r] \& {} \end{tikzcd}}} \deleted{\unsure{ou}}
etc. Comme le foncteur
\[ \left\{\begin{aligned}
X&\longmapsto \deleted{\unsure{\underline{\mathrm{Ch}}}}\mathrm{Ch}_\tau(X)\quad\unsure{(}(=\mathrm{Ob}\,\underline{\mathrm{Ch}}_\tau(X))\\
\mathrm{Cat}&\longrightarrow(\mathrm{Ens})
\end{aligned}\right. \]
est \uline{représenté} par l'objet $I_\tau$ de $X$, \unsure{le sorite|les sorites} des foncteurs représentables nous dit que de telles lois de formation \inserted{de chemins de type $\tau$ \deleted{\unsure{à}} à partir de chemins de type $\tau'$} correspondent \deleted{\unsure{biunivoquement}} \inserted{biunivoquement} \deleted{(à isom. unique près)} aux \deleted{\unsure{fonct}} flèches de $\mathrm{Cat}$
\diagram{\begin{tikzcd}[ampersand replacement=\&] I_{\tau'} \arrow[r,"{\varphi}"] \& I_\tau \end{tikzcd}},
en prenant simplement la ``loi'' correspondante à $\varphi$
\[ c\longmapsto c\circ\varphi\qquad\text{pour }c\colon I_\tau\longrightarrow X, \]
De façon \inserted{plus} précise, \deleted{\unsure{la correspondance}} \inserted{l'application} \deleted{\unsure{la loi}} naturelle
\[ \begin{aligned}
\mathrm{Hom}_{\mathrm{Cat}}(I_{\tau'},I_\tau)&\longrightarrow\mathrm{Hom}(\mathrm{Ch}_{\unsure{\tau}},\mathrm{Ch}_{\tau'})\\
\mathrm{Ob}\,\underline{\mathrm{Hom}}(I_{\tau'},I_\tau)&
\end{aligned} \]
% archive sheet number: 279 (last digit unclear)
\scanpage{281}
\mspage{22}

où le deuxième $\mathrm{Hom}$ est considéré comme un $\mathrm{Hom}$ d'objets dans $\mathrm{Cat}^{\unsure{\times|\vee}}=\underline{\mathrm{Hom}}((\mathrm{Cat}),\mathrm{Ens})$.

Ainsi, dans les deux exemples plus haut, où $\tau=\mathbb{\Delta}_2^+$ et $\tau'=\mathbb{\Delta}_1^+$ ou $\mathbb{\Delta}_2^+$, les foncteurs en sens inverse sont
\diagram{\begin{tikzcd}[ampersand replacement=\&,row sep=small,column sep=small]
\mathbb{\Delta}_2^+\quad\{ \& 0 \arrow[r] \& 1 \arrow[r] \& 2 \& \unsure{/}\\
\mathbb{\Delta}_1^+\quad\{ \& 0 \arrow[u,maps to] \arrow[rr] \&{} \& 1 \arrow[u,maps to] \&{}
\end{tikzcd}}
et
\diagram{\begin{tikzcd}[ampersand replacement=\&,row sep=small,column sep=small]
\mathbb{\Delta}_2^+\quad\{ \&0 \arrow[r] \& 1 \arrow[r,dash] \& 2\\
\mathbb{\Delta}_2^+\quad\{ \&0 \arrow[r] \arrow[u,maps to] \& 1 \arrow[r] \arrow[ur,maps to] \& 2 \arrow[u,maps to]
\end{tikzcd}}.

C'est là un principe général pour ramener toutes les ``opérations \inserted{$\pm$} canoniques sur des chemins \inserted{dans des cat. arbitraires}, et des énoncés les concernant, à des questions où $X$ disparaît, et où on travaille dans les catégories intervalles elles-mêmes.

Un des avantages techniques, c'est que ce sont des catégories fort petites, donnant lieu de plus à une intuition topologique \inserted{sûre} (qui risque de se volatiliser \inserted{par moments} en travaillant \deleted{\unsure{dans}} dans des catégories de chemins), et aussi que dans \unsure{de|des} telles catégories, \deleted{et les} \inserted{donc aussi \unsure{des|les}} catégories de foncteurs d'entre elles, il y a tjrs au plus
% archive sheet number: 280
\scanpage{282}
\mspage{23}

\uline{une seule \unsure{flèche|flèches}} allant d'un objet donné vers un autre. En d'autres termes, dans de telles catégories tous les diagrammes sont commutatifs.

Bien sûr, on voit aussi que si on a une loi de formation \inserted{$L$} de chemins de type $\tau'$ à partir de chemins de type $\tau$ (avec la condition de fonctorialité p.r. à la cat. $X$ où on prend les chemins), cette loi est nécessairement \uline{fonctorielle par rapport aux chemins} : si on \deleted{\unsure{se do|se de}} a un hom
\diagram{\begin{tikzcd}[ampersand replacement=\&]c \arrow[r,"{\unsure{u|v}}"] \&c' \&\text{dans}\quad\underline{\mathrm{Ch}}_\tau(X)\end{tikzcd}}
alors la loi \inserted{$L$} ``se prolonge'' à de \unsure{tels|telles} $u$ pour donner
\diagram{\begin{tikzcd}[ampersand replacement=\&]L(c) \arrow[r,"{L(u)}"] \&L(c')\end{tikzcd}},
avec les relations habituelles
\[ L(\mathrm{id}_c)=\mathrm{id}_{L(c)},\quad L(u'u)=L(u')L(u). \]

\unsure{Cela|Celà} résulte de la description de $L$ comme une opération $L_\varphi$ en termes d'un foncteur bien déterminé
\diagram{\begin{tikzcd}[ampersand replacement=\&] I_{\tau'}\arrow[r,"{\varphi}"] \& I_\tau\end{tikzcd}},
lequel à son tour donne pour toute $X$ un \uline{foncteur}
\[ \left\{\begin{aligned}
\underline{\mathrm{Ch}}_\tau(X)&\longrightarrow\underline{\mathrm{Ch}}_{\tau'}(X)\\
c&\longmapsto c\circ\varphi
\end{aligned}\right. \]
plus précis que simplement une \uline{application} $\mathrm{Ch}_\tau(X)\longrightarrow\mathrm{Ch}_{\tau'}(X)$ entre les ens. d'objets correspondants\unsure{--|.}
% archive sheet number: 281
\scanpage{283}
\mspage{24}

On a le choix \deleted{entre} maintenant entre

a) Une étude plus ou moins systématique des foncteurs entre intervalles. Je préfère attendre pour cela d'avoir besoin de ceci et de cela, et de ne faire que le strict minimum pour les constructions que j'ai en vue.

b) Une description de la composition des chemins.

c) \unsure{Un|Une} formalisme des catégories de chemins \uline{de type non précisé}.

Comme mon objectif est \unsure{avant tout|cependant|maintenant} de décrire de telles catégories $\underline{\mathrm{Ch}}(X)$ \inserted{(sans indice $\tau$ !)}, avec des bonnes propriétés, je commencerai par c, \deleted{et ceci} regarderai ensuite b) comme un test crucial pour l'utilisabilité des notions obtenues, et enfin m'occuperai de a) quand je saurai de quoi j'aurai \unsure{eu|en} besoin.

4. \uline{Catégorie de chemins} $\underline{\mathrm{Ch}}(X)$.

Je veux ``inclure'' \deleted{\unsure{de les|toutes}} les catégories $\underline{\mathrm{Ch}}_\tau(X)$ de chemins de tous les types qu'on veut, dans une $=$ catégorie $\underline{\mathrm{Ch}}(X)$. Rappelons
\[ \underline{\mathrm{Ch}}_\tau(X)=\underline{\mathrm{Hom}}(I_\tau,X), \]
où $I_{\unsure{\tau}}$ est l'\deleted{\unsure{catég}}intervalle (orienté) \unsure{standard|standart} de type $\tau$. \deleted{On va considérer les $\unsure{I}_\tau$ comme des objets}

L'ensemble de tous les types d'intervalles
% archive sheet number: 282
\scanpage{284}
\mspage{26}

La composition des morphismes est la composition des foncteurs entre catégories $I_{\unsure{?|\tau}}$. Ainsi on a un foncteur \uline{pleinement fidèle}$^{*}$ canonique
% slanted
\margin{$*$ faux, cf \unsure{annotation marginale|annotations marginales} page précédente --- on ne prend pas \uline{tous} les foncteurs \unsure{dans|de} $\mathrm{Cat}$ entre $I_\tau$ et $I_{\tau'}$.}
% slanted
\margin{foncteur \uline{réalisation} d'un type (combinatoire) d'intervalle}
\[ \tau\longmapsto I_\tau\colon\quad \underline{\mathrm{Ty}}\xrightarrow{\unsure{\rho|\varphi}}(\mathrm{Cat}), \]
par quoi la catégorie $I_\tau$ \deleted{\unsure{pourrait}} s'identifie à une sous-catégorie pleine de $(\mathrm{Cat})$ (celle des $I_\tau$). Mais je préfère ne pas identifier $\tau\in\mathrm{Ty}$ (donnée combinatoire) avec $I_\tau$ (donnée catégorique).

Soit $X$ une \inserted{(petite ?)} catégorie. On a un contra-foncteur
\[ \left\{\begin{aligned} S&\longmapsto\underline{\mathrm{Hom}}(S,X)\\ (\mathrm{Cat})^{\mathrm{o}}&\mathbin{\text{---}}(\mathrm{Cat})\end{aligned}\right. \]
défini par $X$, d'où en composant avec $\rho$ un contrafoncteur
\diagram{\begin{tikzcd}[ampersand replacement=\&]\underline{\mathrm{Ty}}^{\mathrm{o}} \arrow[r,"{\mathrm{CH}_X}"] \& \mathrm{Cat}\end{tikzcd}}

C'est le contrafoncteur
\[ \left\{\begin{aligned}
\tau&\longmapsto\underline{\mathrm{Hom}}(I_\tau,X)\overset{\mathrm{def}}{=}\underline{\mathrm{Ch}}_\tau(X).\\
&\text{i.e.}\quad \mathrm{CH}_X(\tau)=\underline{\mathrm{Ch}}_\tau(X)
\end{aligned}\right. \]

Mais un contra-foncteur sur une petite catégorie $\underline{\mathrm{Ty}}$ donne lieu à une \uline{catégorie fibrée} (et même scindée) sur celle-ci.
% archive sheet number: 283
\scanpage{285}
\mspage{25}

peut être identifié au monoïde libre \inserted{$\mathrm{Ty}$} à deux générateurs $\Delta^+$ et $\Delta^-$, dont les éléments \deleted{s'obtiennent en faisant} sont les suites finies quelconques (evt vides) d'objets alternativement \unsure{égales|égaux} à $\Delta^+$ et à $\Delta^-$. \deleted{\unsure{Un}}

Si le premier élément est $\Delta^+$, on a un \deleted{\unsure{a}} type de chemin à di\uline{rec}tion initiale directe, rétrograde dans le cas contraire. La décomposition de $I_\tau$ \inserted{ou de $\mathrm{Fl}(I_\tau)$} en intervalles critiques correspond à l'écriture réduite d'un mot comme
\[ {\Delta^+}^{\ell_1}(\Delta^-)^{\ell_2}\ldots(\Delta^{\unsure{?}})^{\ell_{\unsure{r|v}}}
\quad\left(\begin{aligned}\Delta^{?}&=\Delta^+&&\text{si }r\text{ impair}\\\Delta^{?}&=\Delta^-&&\text{si }r\text{ pair}\end{aligned}\right) \]
ou comme
\[ (\Delta^-)^{\ell_1}(\Delta^+)^{\ell_2}\ldots(\Delta^{?})^{\ell_r}
\quad\left(\begin{aligned}\Delta^{?}&=\Delta^-&&\text{si }r\text{ impair}\\
\deleted{\unsure{\delta}}\Delta^{?}&=\Delta^+&&\text{si }r\text{ pair}\end{aligned}\right). \]
Les $\ell_i$ sont les longueurs des intervalles critiques successifs (ou des chaînes structurales successives), leur nombre est la longueur réduite de l'intervalle.

On va considérer $\mathrm{Ty}$ comme l'ens des objets d'une catégorie $\underline{\mathrm{Ty}}$
\[ \left\{\begin{aligned}
\mathrm{Ob}\,\underline{\mathrm{Ty}}&=\mathrm{Ty}=\begin{aligned}[t]
&\text{\deleted{ensemble des} monoïde libre à deux}\\
&\text{générateurs }\inserted{\Delta^+,\Delta^-}\ \text{\unsure{précisés|précédents}.}
\end{aligned}\\
\text{Si }\tau,\tau'&\in\mathrm{Ty},\quad \mathrm{Hom}(\tau,\tau')=\mathrm{Hom}_{\unsure{\mathrm{Cat}}}(I_\tau,I_{\tau'}).
\end{aligned}\right. \]
% slanted
\margin{Définition un peu \unsure{cavalière}, il faut ne prendre que les $\varphi$ de $I_\tau$ dans $I_{\tau'}$ compatibles avec les \deleted{\unsure{objets}} sources et buts --- cf \unsure{plus bas|p. suiv.}.}

\unsure{On notera qu'ici|On notera qu'on a|On peut même préciser}, \deleted{\unsure{non}} seulement un ensemble $\mathrm{Hom}_{\mathrm{Cat}}(I_\tau,I_{\tau'})$, mais même une \uline{catégorie} $\underline{\mathrm{Hom}}(I_\tau,I_{\tau'})$, ce qui revient ici à dire que $\mathrm{Hom}(I_\tau,I_{\tau'})$ (\unsure{i.e.} $\mathrm{Hom}(\tau,\tau')$) peut être vu comme un ens. \uline{ordonné}. Mais pour l'instant nous allons ignorer cette relation d'ordre.
% archive sheet number: 284
\scanpage{286}
\mspage{27}

Je vais \deleted{la} désigner cette catégorie fibrée par $\underline{\mathrm{Ch}}\unsure{\cdot|.}(X)$. (Provisoirement -- si elle s'avère trop grosse, je la \unsure{restreindrai|restraindrai}\ldots). Ainsi on a
\diagram{\begin{tikzcd}[ampersand replacement=\&]\underline{\mathrm{Ch}}(X) \arrow[d,"{\pi_X}"]\\ \underline{\mathrm{Ty}}\end{tikzcd}}.

La fibre de $\underline{\mathrm{Ch}}(X)$ en $\tau\in\mathrm{Ob}\,\underline{\mathrm{Ty}}$ est la catégorie $\underline{\mathrm{Ch}}_\tau(X)$ des chemins de type $\tau$ dans $X$. \deleted{Mais \unsure{qu}} Cela nous dit donc quelles sont les flèches de \deleted{\unsure{$X$}}\inserted{$\underline{\mathrm{Ch}}(X)$} au dessus des flèches identiques de $\underline{\mathrm{Ty}}$. Mais que sont les flèches \inserted{plus} générales ? On se donne une flèche dans $\underline{\mathrm{Ty}}$
\deleted{$I_{\tau'}\longrightarrow I_\tau$}
\[ \varphi\colon\tau'\longrightarrow\tau\qquad\text{i.e.}\quad I_{\tau'}\xrightarrow{\varphi}I_\tau\quad\text{dans }\mathrm{Cat}, \]
quelles sont les flèches de $\underline{\mathrm{Ch}}(X)$ au dessus de $\varphi$ ? Par le sorite général des catégories fibrées, si
\[ \begin{aligned}
c&\in\underline{\mathrm{Ch}}(X)_\tau=\underline{\mathrm{Ch}}_\tau(X)\overset{\mathrm{def}}{=}\underline{\mathrm{Hom}}(I_\tau,X)\\
c'&\in\underline{\mathrm{Ch}}(X)_{\tau'}=\underline{\mathrm{Ch}}_{\tau'}(X)\overset{\mathrm{def}}{=}\underline{\mathrm{Hom}}(I_{\tau'},X),
\end{aligned} \]
on a
\[ \boxed{\mathrm{Hom}_\varphi(\unsure{c'|\tau'},c)\simeq\mathrm{Hom}\left(c',\underbrace{\varphi^*(c)}_{\substack{\Vert\\c\circ\varphi}}\right)} \]
Donc \unsure{!|:} les \deleted{flèches \unsure{dans}} homomorphismes d'un
% archive sheet number: 285
\scanpage{287}
\mspage{28}

chemin \deleted{\unsure{$v$}} $c'$ de type $\tau'$ dans $X$, dans le chemin $c$ de type $\tau$, ``au dessus'' \unsure{du|de} \deleted{\unsure{le}} morphisme $\varphi$ de $\tau'$ dans $\tau$ i.e. du foncteur
\[ \varphi\colon I_{\tau'}\longrightarrow I_\tau, \]
est un hom \inserted{$u$} de $c'$ dans $c\circ\varphi$, i.e.
\diagram{\begin{tikzcd}[ampersand replacement=\&,row sep=large,column sep=large]
I_{\tau'} \arrow[r,"{\varphi}"] \arrow[dr,"{c'}"',""{name=leftpath,coordinate,pos=.60}] \& I_\tau \arrow[d,"c",""{name=rightpath,coordinate,pos=.55}]\\
{} \& X \arrow[from=leftpath,to=rightpath,bend left=12,"u"{circle,draw,inner sep=1pt}]
\end{tikzcd}}

La composition des flèches est évidente :
\diagram{\begin{tikzcd}[ampersand replacement=\&,row sep=large,column sep=large]
I_{\tau''} \arrow[r,"{\varphi'}"] \arrow[drr,bend right=27,"{c''}"',""{name=outerpath,coordinate,pos=.62}] \&
I_{\tau'} \arrow[r,"{\varphi}"] \arrow[dr,"{c'}"',""{name=midpath,coordinate,pos=.57}] \&
I_\tau \arrow[d,"c",""{name=innerpath,coordinate,pos=.50}]\\
{} \&{} \& X \arrow[from=outerpath,to=midpath,bend left=14,"{u'}"] \arrow[from=midpath,to=innerpath,bend left=14,"u"]
\end{tikzcd}}
on a :
\[ \begin{aligned}
u&\colon c'\longrightarrow c\circ\varphi,\quad\text{d'où}\quad c'\circ\varphi'\xrightarrow{u*\varphi'}\underset{\substack{\Vert\\c\circ(\varphi\varphi')}}{(c\circ\varphi)\unsure{\circ}\varphi'}\\
u'&\colon c''\mathbin{\text{---}}c'\circ\varphi'
\end{aligned} \]
et $u''=(u*\varphi')\circ u'$ \unsure{donné|donc}
\diagram{\begin{tikzcd}[ampersand replacement=\&,column sep=small] c''\arrow[r,"{u'}"] \arrow[rr,bend right=22,"{u''}"'{circle,draw,inner sep=1pt}] \&c'\circ\varphi'\arrow[r,"{u*\varphi'}"] \&c\circ(\varphi\varphi') \end{tikzcd}}

Mais a-t-on des foncteurs source et but
\diagram{\begin{tikzcd}[ampersand replacement=\&]\underline{\mathrm{Ch}}(X) \arrow[r,shift left=.6ex,"{\unsure{\sigma|\alpha}}"] \arrow[r,shift right=.6ex,"{\beta}"'] \& X \end{tikzcd}}
\[ \begin{aligned} \unsure{c|\tau}&\longmapsto\mathrm{source}(c)\\ c&\longmapsto\mathrm{but}(c)\end{aligned}\quad ? \]

Il faudrait donc qu'un hom $u$ comme dessus
\[ u\colon c'\deleted{\unsure{\circ\varphi}}\mathbin{\text{---}}c\circ\varphi \]
définisse
\[ \begin{aligned} \sigma(u)&\colon\mathrm{source}(c')\mathbin{\text{---}}\mathrm{source}(c)\\ \beta(u)&\colon\deleted{\unsure{\mathrm{s}}}\mathrm{but}(c')\mathrel{\unsure{\text{---}|\longrightarrow}}\mathrm{but}(c)\end{aligned} \]
% archive sheet number: 286
\scanpage{288}
\mspage{29}

Mais comme \deleted{\unsure{la}}
\[ \left\{\begin{aligned}\unsure{\sigma|\alpha}(c)&=c(0_{\deleted{\unsure{?}}\tau}),\quad \beta(c)=c(1_\tau)\\
&\text{où }0_\tau,1_\tau\text{ sont les objets source et but}\\
&\text{de l'intervalle orienté \unsure{standard|standart} }I_\tau
\end{aligned}\right. \]
et de $=$ pour $\tau'$, \unsure{$u$|il} \deleted{\unsure{est|n'est}} \unsure{induira|n'induira} bien
\[ \text{\deleted{\unsure{source}}}\quad\left\{\begin{aligned}
\sigma(c')&\longrightarrow c(\varphi(0_{\tau'})\overset{?}{=}\sigma(c)\\[-3pt]
\Vert &\hspace{10em}\Vert\\
c'(0_{\tau'})&\hspace{9em}c(0_\tau)
\end{aligned}\right., \]
et idem pour $\beta(c')$ et $\beta(c)$, que si on impose
\[ \boxed{\left\{\begin{aligned}\varphi(0_{\tau'})&=0_\tau\\ \varphi(1_{\tau'})&=1_\tau\end{aligned}\right.} \]
i.e. que le foncteur envisagé $\varphi$ soit \uline{compatible avec les sources et les buts}. C'est ça \inserted{visiblement} qu'il fallait imposer dans \deleted{\unsure{$\varphi$}} la définition de
\[ \mathrm{Hom}_{\underline{\mathrm{Ty}}}(\tau',\tau)=\underline{\mathrm{Hom}}^{\tikz[baseline=(t.base)]{\node[draw,ellipse,inner sep=.6pt] (t) {\scriptsize int};}}(I_{\tau'},I_\tau)\quad - \]
il faut prendre les foncteurs \inserted{$\varphi$} respectant les structures d'intervalle, i.e. rendant commutatif \unsure{le|les} diagramme
\diagram{\begin{tikzcd}[ampersand replacement=\&,column sep=small]
I_{\tau'} \arrow[rr] \&{} \& I_\tau\\
\deleted{\unsure{\varepsilon}} \&\varepsilon=\{0,1\} \arrow[ul,hook,"{i_{\tau'}}"'] \arrow[ur,hook,"{i_\tau}"']\&{}
\end{tikzcd}}
% slanted
\margin{En fait, il faut sans doute \unsure{restreindre|contraindre} encore plus les morphismes $\varphi$, en exigeant qu'ils soient \uline{croissants} sur les objets --- cf page 31.}

définissant l'orientation de $I_{\tau'}$ et de $I_\tau$. Cette rectification faite, j'ai bon espoir que la catégorie $\underline{\mathrm{Ch}}(X)$ qu'on vient de construire est ``la bonne''. (On pourrait désigner par $\mathrm{CH}(X)$ la catégorie plus grande, sur la catégorie $\underline{\unsure{\mathrm{Typ}|\mathrm{Ty}_p}}$ \deleted{plus} ayant ${=}^{\mathrm s}$ objets que $\underline{\mathrm{Ty}}$, mais plus de morphismes, \unsure{à savoir}
% archive sheet number: 287
\scanpage{289}
\mspage{30}

\[
\begin{aligned}
\mathrm{Hom}_{\underline{\mathrm{Typ}}}(\tau,\tau')&=\mathrm{Hom}_{\mathrm{Cat}}(I_\tau,I_{\tau'})\\
\mathrm{Hom}_{\mathrm{Ty}}(\tau,\tau')&=\mathrm{Hom}^{\mathrm{int}}_{\mathrm{Cat}}(I_\tau,I_{\tau'})
\end{aligned}
\]
(hom. respectant la structure d'intervalle)

On aura
\[
\underline{\mathrm{C}\deleted{\unsure{\mathrm{H}}}\mathrm{h}}(X)\simeq \underline{\mathrm{CH}}(X)\mathbin{\times}_{\underline{\mathrm{Typ}}}\underline{\mathrm{Ty}},
\]
i.e. $\underline{\mathrm{Ch}}(X)$ est la \uline{restriction} de la catégorie fibrée $\underline{\mathrm{CH}}(X)$ sur $\underline{\mathrm{Typ}}$, à la sous-catégorie (non pleine) $\underline{\mathrm{Ty}}$ de $\underline{\mathrm{Typ}}$.

\uline{NB} Les catégories fibrées $\underline{\mathrm{CH}}(X),\underline{\mathrm{Ch}}(X)$ sur $\underline{\mathrm{Typ}},\underline{\mathrm{Ty}}$ dépendent fonctoriellement de $X$, variant dans $(\mathrm{Cat})$, comme il se doit. Comme de juste, on s'y intéresse surtout quand $X$ \uline{n'est pas} dans $(\mathrm{Cat})$, i.e. quand $X$ est une ``grosse'' catégorie (telle que $(\mathrm{Cat})$ elle-même). Je me dispense \deleted{\unsure{de}}d'expliciter ici les questions pertinentes d'Univers.

Deux tests essentiels, pour s'assurer si la notion est bonne

a) La théorie des $\unsure{M|\mathcal M}\Sigma^{-1}$

b) La composition des chemins) doit bien marcher.
% archive sheet number: 288
\scanpage{290}
\mspage{31}

On a le diagramme de foncteurs canoniques
\diagram{\begin{tikzcd}[ampersand replacement=\&,column sep=large]
\underline{\mathrm{Ty}} \arrow[r,hook] \arrow[d,hook] \& {}_{\varepsilon}\backslash\mathrm{Cat} \arrow[d]\\
\underline{\mathrm{Typ}} \arrow[r,arrows={Hooks[harpoon]-}] \& \mathrm{Cat}\mathpunct{,}
\end{tikzcd}}
où les flèches horizontales sont pl. fidèles.

5. \uline{Propriétés} \deleted{de $\underline{\mathrm{Ty}}$, \unsure{et}} \uline{des morphismes d'intervalles}

Je me suis plus ou moins convaincu que pour se prêter à des calculs et à des raisonnements simples, il faut encore plus restreindre les morphismes dans la \inserted{définition de la} catégorie $\underline{\mathrm{Ty}}$. Je \deleted{\unsure{vais}} exige\deleted{\unsure{r}} que les morphismes $\tau\to\tau'$ soient des foncteurs
\[
\unsure{w|u}\colon I_\tau\mathrel{\unsure{\text{---}|\longrightarrow}}I_{\tau'},
\]
qui sont compatibles avec la structure d'\uline{intervalles} \uline{orientés} de $I_\tau,I_{\tau'}$, i.e. \inserted{qu'ils} transforment source en source, but en but, mais de plus qu'ils soient \uline{croissants} sur les objets, (pour la relation d'ordre de progression des chemins).
C'est une restriction qui géométriquement me parait assez naturelle. (Elle est similaire à la restriction analogue, dans la définition pré-catégorique des morphismes entre $\mathbb{\Delta}_{\unsure{i|1}}\ldots$) C'est cette hypothèse qui permet de faire une théorie de structure assez simple des \deleted{foncteurs} morphismes d'intervalles. Mais pour l'instant je vais distinguer ``morphismes'' et ``morphismes croissants''
% archive sheet number: 289
